% Official published PDF pp.29:2--4; original numbering retained in headings.
\section*{Partial commutative sums: original Definitions 1--4, 10 and 12}
If $A$ is a countable set and $M$ is a set, an $A$-indexed family of elements of $M$ is a function $\vec x:A\to M$, written $\vec x=(x_a)_{a\in A}$. We consider a set $M$ with a partial function $\Sigma$ called the sum from indexed families on $M$ to $M$. An indexed family $(x_a)_{a\in A}$ is summable if it is in the domain of $\Sigma$, and we write its image as $\sum_{a\in A}x_a$.
\begin{originaldef}[1]
For any expression $e$ and $e'$ involving partial sums, we write $e\dimp e'$ (or $e'\mathrel{\leftarrow\joinrel\relbar}e$) if whenever $e$ is defined, then $e'$ is also defined and $e=e'$. We write $e\dequiv e'$ if $e\dimp e'$ and $e'\dimp e$.
\end{originaldef}
\begin{originaldef}[2, original reference 22]
The tuple $(M,\Sigma)$ is a partial commutative monoid (PCM) if $M$ is non-empty and if the sum $\Sigma$ satisfies the following axioms.
\begin{description}
\item[(Unary)] Unary sum axiom: every singleton family $(x)$ is summable with sum $x$.
\item[(WPA)] Weak partition associativity axiom: Let $(x_a)_{a\in A}$ be an indexed family, $I$ countable, and $\{A_i\}_{i\in I}$ a partition of $A$ (the $A_i$'s can be empty). Then
\[\sum_{a\in A}x_a\dimp\sum_{i\in I}\left(\sum_{a\in A_i}x_a\right).\]
\end{description}
A PCM is strong if it satisfies the partition associativity axiom (PA), where the implication $\dimp$ of (WPA) is replaced by an equivalence $\dequiv$.
\end{originaldef}
Observe that the neutral element $0$ is not given as part of the data, but can be defined from the axioms. Indeed, the sum over the empty family is always defined by Lemma 3.
\begin{originallemma}[3] If $A\subseteq B$ and $(x_b)_{b\in B}$ is summable, then $(x_a)_{a\in A}$ is summable.\end{originallemma}
\begin{originaldef}[4] Every PCM has a canonical preorder (it is not antisymmetric in general) given by $x\leq y$ if there exists $z$ such that $x+z=y$.\end{originaldef}
\begin{originaldef}[10]
A Partial Commutative Rig (PCR) is the data of $(R,\Sigma,1,\cdot)$ such that $(R,\Sigma)$ is a PCM, $(R,1,\cdot)$ is a commutative monoid, and such that for all indexed families $(x_a)_{a\in A}$ and $(y_b)_{b\in B}$ of $R$,
\[\left(\sum_{a\in A}x_a\right)\cdot\left(\sum_{b\in B}y_b\right)\dimp\sum_{a\in A,b\in B}(x_a\cdot y_b).\]
A PCR is strong if the underlying PCM is strong.
\end{originaldef}
\begin{originaldef}[12, original reference 22]
Let $L$ be a locally small category such that for all objects $X,Y$, $L(X,Y)$ is a PCM. The category $L$ is a PCM category if for all $(f_a\in L(X,Y))_{a\in A}$ and $(g_b\in L(Y,Z))_{b\in B}$,
\[\left(\sum_{b\in B}g_b\right)\circ\left(\sum_{a\in A}f_a\right)\dimp\sum_{(a,b)\in A\times B}(g_b\circ f_a).\]
\end{originaldef}
% Official published PDF pp.29:5--7 and29:9--10.
\section*{Web models and predual characterization: original Part I}
We define the Kronecker symbol: $\delta_{ab}=0$ if $a\ne b$, and $\delta_{aa}=1$. We define for all $a\in |X|$ a vector $e_a\in R^{|X|}$ by $(e_a)_{a'}=\delta_{aa'}$. For all set $|X|$ (called the web) and all PCR $R$, we define the following partial maps on vectors:
\[\scalar{x}{y}_R=\sum_{a\in|X|}x_ay_a,\qquad (t\cdot s)_{a,c}=\sum_{b\in|Y|}s_{a,b}\cdot t_{b,c},\qquad(s\cdot x)_b=\sum_{a\in|X|}s_{a,b}x_a.\]
Let $L$ be a model of linear logic, more precisely a Seely category. We write the composition of $f\in L(X,Y)$ with $g\in L(Y,Z)$ as $g\circ f$. We write the symmetric monoidal structure as $(\otimes,\lambda,\rho,\alpha,\gamma)$ where $\lambda_X:1\otimes X\to X$ and $\rho_X:X\otimes1\to X$ are the unitors, $\alpha_{X,Y,Z}:(X\otimes Y)\otimes Z\to X\otimes(Y\otimes Z)$ is the associator, and $\gamma_{X,Y}:X\otimes Y\to Y\otimes X$ is the symmetry. We write $X\lolli Y$ the internal hom, the currying of $f:X\otimes Y\to Z$ is $\cur(f):X\to(Y\lolli Z)$, and the evaluation is $\mathrm{ev}:(X\lolli Y)\otimes X\to Y$.
First we assume that $R_+$ is a strong PCR, and we consider $B\subseteq R_+$. We assume that $B$ is downward closed: if $x\leq y$ and $y\in B$ then $x\in B$. We define an orthogonality relation $x\perp y$ on vectors and the orthogonal of a set $F\subseteq R_+^{|X|}$ as:
\[x\perp y\ \text{if }\scalar{x}{y}_R\text{ is defined and belongs to }B,\qquad F^\perp=\{x'\in R_+^{|X|}\mid x\perp x'\}.\]
This construction has the usual properties of an orthogonality construction: $F\subseteq F^{\perp\perp}$, $(F\subseteq G)\Rightarrow(G^\perp\subseteq F^\perp)$, and $F^{\perp\perp\perp}=F^\perp$.
\begin{originaldef}[14] A set of vectors $F\subseteq R_+^{|X|}$ is a covering if for all $a\in|X|$, there exists $x\in F$ such that its $a$ component $x_a$ is invertible.\end{originaldef}
\begin{originaldef}[15] An $R_+$-space is a pair $X=(|X|,AX)$ where $|X|$ is a set and $AX\subseteq R_+^{|X|}$ is a set of vectors such that $AX^{\perp\perp}=AX$, and such that both $AX$ and $AX^\perp$ are coverings. The dual of an $R_+$-space $X$ is defined as the $R_+$-space $X^\perp=(|X|,AX^\perp)$.\end{originaldef}
\begin{originallemma}[16, Downward closure] For all $x,y\in R_+^{|X|}$ such that $x\leq y$, $y\in AX\Rightarrow x\in AX$.\end{originallemma}
For all vectors $x\in R_+^{|X|}$ and $y\in R_+^{|Y|}$, define $x\otimes y\in R_+^{|X|\times|Y|}$ by $(x\otimes y)_{a,b}=x_ay_b$. For all sets $F\subseteq R_+^{|X|}$ and $G\subseteq R_+^{|Y|}$, we set $F\otimes G=\{x\otimes y\mid x\in F,y\in G\}$. The tensor $X\otimes Y$ of two $R_+$-spaces $X$ and $Y$ is defined as
\[|X\otimes Y|=|X|\times|Y|,\qquad A(X\otimes Y)=(AX\otimes AY)^{\perp\perp}.\]
We now define the $R_+$-space of linear maps between $R_+$-spaces by duality: $X\lolli Y=(X\otimes Y^\perp)^\perp$.
\paragraph{Remark 18.} $A(X\lolli Y)=A(X\otimes Y^\perp)^\perp=(AX\otimes AY^\perp)^{\perp\perp\perp}=(AX\otimes AY^\perp)^\perp$. Thus, a matrix $s\in R^{|X\lolli Y|}$ belongs to $A(X\lolli Y)$ if and only if $s\perp(x\otimes y')$ for all $x\in AX$, $y'\in AY^\perp$.
\begin{originalthm}[19] For all $s\in A(X\lolli Y)$ and $t\in A(Y\lolli Z)$, $t\cdot s$ is well-defined and belongs to $A(X\lolli Z)$. Furthermore, the identity matrix $\id_X\in R_+^{|X\lolli X|}$ belongs to $A(X\lolli X)$.\end{originalthm}
Thus, we can define the category $\Web_{R_+}$ whose objects are the $R_+$-spaces, whose morphisms are $\Web_{R_+}(X,Y)=A(X\lolli Y)$, and whose composition is given by matrix multiplication.
\begin{originallemma}[20] For all $s\in R_+^{|X|\times|Y|}$, $x\in R_+^{|X|}$ and $y\in R_+^{|Y|}$,
\[\scalar{s\cdot x}{y}\dimp\scalar{s}{x\otimes y}.\]
Furthermore, if $s\cdot x$ is defined, then $\scalar{s\cdot x}{y}\dequiv\scalar{s}{x\otimes y}$.
\[\scalar{s^\perp\cdot y}{x}\dimp\scalar{s}{x\otimes y}.\]
Furthermore, if $s^\perp\cdot y$ is defined, then $\scalar{s^\perp\cdot y}{x}\dequiv\scalar{s}{x\otimes y}$. Where $s^\perp\in R_+^{|Y|\times|X|}$ is the transpose of $s$, $(s^\perp)_{b,a}=s_{a,b}$.
\end{originallemma}
\begin{originallemma}[21, Covering principle] Let $Q$ be a covering of $R_+^{|Y|}$. Let $x\in R_+^{|X|}$ and $s\in R_+^{|X\lolli Y|}$ such that $\scalar{s}{x\otimes y}$ is well-defined for all $y\in Q$. Then $s\cdot x$ is well-defined. Similarly, let $P$ be a covering of $R_+^{|X|}$. Let $y\in R_+^{|Y|}$ and $s\in R_+^{|X\lolli Y|}$ such that $\scalar{s}{x\otimes y}$ is well-defined for all $x\in P$. Then $s^\perp\cdot y$ is well-defined.\end{originallemma}
\begin{originalprop}[22] Let $s\in R_+^{|X\lolli Y|}$. The following are equivalent:
\begin{enumerate}\item $s\in A(X\lolli Y)$;\item $s^\perp\in A(Y^\perp\lolli X^\perp)$;\item $\forall x\in AX$, $s\cdot x$ is defined and belongs to $AY$;\item $\forall y'\in AY^\perp$, $s^\perp\cdot y'$ is defined and belongs to $AX^\perp$.\end{enumerate}
\end{originalprop}
\begin{originalthm}[23, Predual characterization of morphisms]
Let $X$ and $Y$ be two $R_+$-spaces. Let $P\subseteq R_+^{|X|}$ and $Q\subseteq R_+^{|Y|}$ be two coverings such that $P^{\perp\perp}=AX$ and $Q^\perp=AY$. Then for all $s\in R_+^{|X\lolli Y|}$, the following are equivalent.
\begin{enumerate}\item $s\in A(X\lolli Y)$;\item $s\in(P\otimes Q)^\perp$;\item for all $x\in P$, $s\cdot x$ is defined and belongs to $AY$;\item for all $y'\in Q$, $s^\perp\cdot y'$ is defined and belongs to $AX^\perp$.\end{enumerate}
\end{originalthm}
\begin{originalcor}[24] Let $P\subseteq R_+^{|X|}$ and $Q\subseteq R_+^{|Y|}$ be two coverings such that $P^{\perp\perp}=AX$ and $Q^{\perp\perp}=AY$. Then $A(X\otimes Y)=(P\otimes Q)^{\perp\perp}$.\end{originalcor}
The following result immediately follows from Theorem 23, using that $X\otimes Y=(X\lolli Y^\perp)^\perp$.
\begin{originaldef}[26] A PCR $R$ is absolute if there exists a strong PCR $R_+$ and a function $|.|:R\to R_+$ (called the absolute value) such that
\begin{itemize}\item $|x|=0$ if and only if $x=0$;\item $|.|$ is a monoid morphism for the multiplicative part: $|xy|=|x||y|$ and $|1|=1$;\item Triangle inequality: $(x_a)_{a\in A}$ is summable in $R$ if and only if $(|x_a|)_{a\in A}$ is summable in $R_+$, and $|\sum_{a\in A}x_a|\leq\sum_{a\in A}|x_a|$.\end{itemize}
\end{originaldef}
The absolute value $|.|$ extends to a function $\abs:R^{|X|}\to R_+^{|X|}$ for all web $|X|$, defined by $\abs(x)_a=|x_a|$ for all $a\in|X|$.
\begin{originaldef}[28] For all $R_+$-space $X$, let $\Rmod{X}=\{x\in R^{|X|}\mid\abs(x)\in AX\}$. We define the category $\Web_R$ whose objects are the $R_+$-spaces, whose morphisms from $X$ to $Y$ are the matrices in $\Rmod{X\lolli Y}$, and whose composition is matrix multiplication.\end{originaldef}
For all $s\in\Rmod{X\lolli Y}$ and $t\in\Rmod{Y\lolli Z}$, we have that $\abs(t\cdot s)\dlequiv\abs(t)\cdot\abs(s)\in A(X\lolli Z)$, so $t\cdot s$ is well-defined and belongs to $\Rmod{X\lolli Z}$ by downward closure of $A(X\lolli Z)$ (Lemma 16).
\begin{originaldef}[30] Let $X$ be an $R_+$-space. A family $(x(i))_{i\in I}$ of elements of $AX$ is summable if their pointwise sum $(\sum_{i\in I}x(i))_a=\sum_{i\in I}x(i)_a$ is defined and belong to $AX$.\end{originaldef}
\begin{originalprop}[31] Let $X$ be an $R_+$-space. Then $AX$ is a strong PCM.\end{originalprop}
\begin{proof}It immediately follows from (PA) on $R_+$ that for all family $(x(i))_{i\in I}$ of $R_+^{|X|}$ and all partition $\{I_k\}_{k\in K}$ of $I$, $\sum_{i\in I}x(i)\dequiv\sum_{k\in K}\sum_{i\in I_k}x(i)$. To conclude, we still need to ensure that whenever $(x(i))_{i\in I}$ is summable in $AX$, then for all $k\in K$, $(x(i))_{i\in I_k}$ is summable in $AX$. This is an immediate consequence of the downward closure of $AX$ (Lemma 16).\end{proof}
Similarly to morphisms, summability can be characterized in terms of predual. The proof is very similar to the proof of Lemma 21 and of Theorem 23.
\begin{originalprop}[33] Let $P$ be a covering of $R_+^{|X|}$ such that $P^\perp=AX$. Then a family $(x(i))_{i\in I}$ of $AX$ is summable if and only if for all $x'\in P$, $\scalar{\sum_{i\in I}x(i)}{x'}$ is defined and is in $B$.\end{originalprop}
\begin{originaldef}[34] Let $X$ be an $R_+$-space. A family $(x(i))_{i\in I}$ of elements of $\Rmod{X}$ is summable if $(\abs(x(i)))_{i\in I}$ is summable in $AX$ (in the sense of Definition 30).\end{originaldef}
\begin{originalprop}[35] The pointwise sum defined above is a PCM on $\Rmod{X}$.\end{originalprop}
\begin{proof} Let $(x(i))_{i\in I}$ be summable in $\Rmod{X}$, and $\{I_k\}_{k\in K}$ be a partition of $I$. For all $k\in K$, $(\abs(x(i)))_{i\in I_k}$ is summable in $AX$ by Lemma 3. Thus, $(x(i))_{i\in I_k}$ is summable in $\Rmod{X}$:
\[\sum_{k\in K}\abs\left(\sum_{i\in I_k}x(i)\right)\mathrel{\stackrel{1}{\dleleft}}\sum_{k\in K}\sum_{i\in I_k}\abs(x(i))\mathrel{\stackrel{2}{\dequiv}}\sum_{a\in A}\abs(x_a)\in AX\]
where (1) holds by triangle inequality and (2) by (PA) on the strong PCM $AX$. We conclude that $(\sum_{i\in I_k}x(i))_{k\in K}$ is summable in $\Rmod{X}$ by downward closure of $AX$.\end{proof}
\begin{originaldef}[37] A summability structure on a category $L$ is a tuple $(S,(\pi_i)_{i\in\mathbb N},\sigma,0)$ where
\begin{description}
\item[(S-fun)] $S$ is a functor on $L$, and $\pi_i,\sigma:S\Rightarrow\mathrm{Id}$ are natural transformations called respectively projections and sum, and $0$ is a family of zero morphisms $0_{X,Y}\in L(X,Y)$.
\item[(S-proj)] The $\pi_i$ are jointly monic: $\forall f,g\in L(X,SY)$, $(\forall i\in\mathbb N,\pi_i\circ f=\pi_i\circ g)\Rightarrow f=g$.
\item[(S-sum)] For all objects $X,Y$, there exists a PCM structure on $L(X,Y)$ such that
\begin{enumerate}\item[(S-sum-1)] A family $(f_{i_1,\ldots,i_n})_{i_1,\ldots,i_n\in\mathbb N}$ over $L(X,Y)$ is summable if and only if there exists $f\in L(X,S^nY)$ such that $\pi_{i_1}\circ\cdots\circ\pi_{i_n}\circ f=f_{i_1,\ldots,i_n}$.
\item[(S-sum-2)] If $(f_{i_1,\ldots,i_n})_{i_1,\ldots,i_n\in\mathbb N}$ is summable then $\sum_{i_1,\ldots,i_n\in\mathbb N}f_{i_1,\ldots,i_n}=\underbrace{\sigma\circ\cdots\circ\sigma}_{n}\circ f$.
\item[(S-sum-3)] The sum over the empty family is equal to the zero morphism $0$.
\end{enumerate}\end{description}\end{originaldef}

% Official published PDF pp.29:14--15.
\section*{Representability: original Definitions 49--50 and Proposition 52}
We assume in this section that $L$ is a symmetric monoidal closed category, with zero morphisms $0$ that are also absorbing for the monoidal product: $\id_X\otimes0=0$ and $0\otimes\id_Y=0$.
\begin{originaldef}[49] A summation object in $L$ is an object $D$ equipped with a family of jointly epic morphisms $(e_i:1\to D)_{i\in\mathbb N}$ called the basis, and $\Delta:1\to D$ called the diagonal.\end{originaldef}
The internal hom $1\lolli X$ can be identified with $X$, and then $\mathrm{ev}_{1,X}=\rho_X:X\otimes1\to X$ and for all $f:X\otimes1\to Y$, $\cur(f)=f\circ\rho_X^{-1}$. We define natural transformations $\pi_i=(e_i\lolli X):(D\lolli X)\to X$ and $\sigma=(\Delta\lolli X):(D\lolli X)\to X$. More explicitly,
\[\pi_i=\mathrm{ev}_{D,X}\circ((D\lolli X)\otimes e_i)\circ\rho^{-1},\qquad\sigma=\mathrm{ev}_{D,X}\circ((D\lolli X)\otimes\Delta)\circ\rho^{-1}.\tag{3}\]
\begin{originaldef}[50] The category $L$ is a representable PCM category if it has a summation object $D$ such that $(D\lolli\_,(\pi_i)_{i\in\mathbb N},\sigma,0)$ is a summability structure.\end{originaldef}
We now provide an alternative characterization of representable PCM categories, which can be obtained by ``uncurrying'' the definition of summability structures. By standard naturality equations on the currying, we get that for all $f:X\otimes D\to Y$ and $g:X\otimes D^{\otimes n}\to Y$,
\begin{align*}\pi_i\circ\cur(f)&=f\circ(X\otimes e_i)\circ\rho^{-1},&\sigma\circ\cur(f)&=f\circ(X\otimes\Delta)\circ\rho^{-1},\\\pi_{i_1}\circ\cdots\circ\pi_{i_n}\circ\cur^n(g)&=g\circ(X\otimes e_{i_n}\otimes\cdots\otimes e_{i_1})\circ(\rho^n)^{-1},\\\sigma\circ\cdots\circ\sigma\circ\cur^n(g)&=g\circ(X\otimes\Delta\otimes\cdots\circ\Delta)\circ(\rho^n)^{-1}.\tag{4}\end{align*}
\begin{originallemma}[51] For all object $D$ and every family of morphisms $(e_i:1\to D)_{i\in\mathbb N}$, define $\pi_i$ as in Equation (3). Then, the following assertions are equivalent:
\begin{enumerate}\item the morphisms $e_i$ are jointly epic;\item for all object $X$, the morphisms $X\otimes e_i:X\otimes1\to X\otimes D$ are jointly epic;\item the morphisms $\pi_i$ are jointly monic.\end{enumerate}\end{originallemma}
\begin{proof}(1) $\Rightarrow$ (2). Let $f,g:X\otimes D\to Y$ be such that $f\circ(X\otimes e_i)=g\circ(X\otimes e_i)$ for all $i\in\mathbb N$. By naturality of $\gamma$, we get $f\circ\gamma\circ(e_i\otimes X)=g\circ\gamma\circ(e_i\otimes X)$ and hence $\cur(f\circ\gamma)\circ e_i=\cur(g\circ\gamma)\circ e_i$ so that $\cur(f\circ\gamma)=\cur(g\circ\gamma)$ and hence $f=g$. So the $X\otimes e_i$ are jointly epic. Conversely, (2) $\Rightarrow$ (1), taking $X=1$. Finally, (2) $\Leftrightarrow$ (3) by Equation (4) and bijectivity of $\cur$.\end{proof}
Thus, the only condition that does not hold de facto when there is a summation object is (S-sum). This condition can be uncurried by Equation (4) as follows.

\begin{originalprop}[52] The category $L$ is a representable PCM category if and only if it has a summation object $D$ such that for all objects $X,Y$, $L(X,Y)$ has a PCM structure in which:
\begin{enumerate}\item A family $(f_{i_1,\ldots,i_n})_{i_n,\ldots,i_n\in\mathbb N}$ of morphisms $X\to Y$ is summable if and only if there exists $h:X\otimes D^{\otimes n}\to Y$ such that $h\circ(X\otimes e_{i_n}\otimes\cdots\otimes e_{i_1})\circ(\rho_X^n)^{-1}=h_{i_1,\ldots,i_n}$.
\item Then, $\sum_{i_1,\ldots,i_n\in\mathbb N}h_{i_1,\ldots,i_n}=h\circ(X\otimes\Delta\otimes\cdots\otimes\Delta)\circ(\rho_X^n)^{-1}$.
\item The sum over the empty family is equal to the zero morphism $0$.
\end{enumerate}\end{originalprop}
% MAIN STATEMENT BEGIN -- official PDF p.29:16.
\section*{Original Section 7: Taylor expansion in web models}
We prove in this section that our web models are representable PCM categories, and feature a Taylor expansion. For all set $A$ and all PCR $R$, define $\Delta_A\in R^A$ by $(\Delta_A)_a=1$. We then define $l^1(A)\subseteq R^A$ and $l^\infty(A)\subseteq R^A$ by $l^1(A)=\{\Delta_A\}^\perp$ and $l^\infty(A)=\{\Delta_A\}^{\perp\perp}$. In particular, $(x_a)_{a\in A}\in l^1(A)$ if and only if $(x_a)_{a\in A}$ is summable and $\sum_{a\in A}x_a\in B$. Then, we set $D=(\mathbb N,l^\infty(\mathbb N))$.
The object $D$ is a summation object, with diagonal $\Delta=\Delta_{\mathbb N}$ and basis $e_i\in R^{|D|}$ defined by $(e_a)_{a'}=\delta_{aa'}$ (they belong to $AD$ by downward closure, since $e_a\leq\Delta$). We now prove that the summability in $\Web_R$ described in Section 4 is representable.
\begin{originallemma}[58] For all $k\in\mathbb N$, $A(D^{\otimes k})=l^\infty(\mathbb N^k)$.\end{originallemma}
\begin{proof}$A(D^{\otimes k})=(\{\Delta\}^{\perp\perp}\otimes\cdots\otimes\{\Delta\}^{\perp\perp})^{\perp\perp}\mathrel{\stackrel{1}{=}}(\{\Delta\}\otimes\cdots\otimes\{\Delta\})^{\perp\perp}=\{\Delta_{\mathbb N^k}\}^{\perp\perp}=l^\infty(\mathbb N^k)$ where (1) is obtained by iterating Corollary 24.\end{proof}
\begin{originalthm}[59] The sum in $\Web_{R_+}$ and $\Web_R$ is representable.\end{originalthm}
% MAIN PROOF BEGIN -- official PDF pp.29:16--17.
\begin{proof}
It suffices to do the proof for $\Web_R$ ($\Web_{R_+}$ is a special instance of $\Web_R$ in which $|.|$ is the identity function). We prove that $D$ satisfies the conditions of Proposition 52. Let $(f(i_1,\ldots,i_k))_{i_1,\ldots,i_k\in\mathbb N}$ be a family of $\Rmod{X\lolli Y}$. Define $h\in R^{|X\otimes D\lolli Y|}$ by $h_{a,i_k,\ldots,i_1,b}=f(i_1,\ldots,i_k)_{a,b}$ so that
\[f(i_1,\ldots,i_k)=h\cdot(X\otimes e_{i_k}\otimes\cdots\otimes e_{i_1})\cdot(\rho^k)^{-1}\]
and
\[\sum_{i_1,\ldots,i_k\in\mathbb N}f(i_1,\ldots,i_k)\dequiv h\cdot(X\otimes\Delta\otimes\cdots\otimes\Delta)\cdot(\rho^k)^{-1}\dequiv h\cdot(X\otimes\Delta_{\mathbb N^k})\cdot(\rho^k)^{-1}.\]
It suffices to prove that $(f(i_1,\ldots,i_k))_{i_1,\ldots,i_k\in\mathbb N}$ is summable in $\Rmod{X\lolli Y}$ if and only if $h\in\Rmod{X\otimes D^{\otimes k}\lolli Y}$ to conclude. We proceed by the following chain of equivalences.
\begin{align*}
&(f_{i_1,\ldots,i_k})_{i_1,\ldots,i_k\in\mathbb N}\text{ is summable in }\Rmod{X\lolli Y}\\
\iff{}&(\abs(f)_{i_1,\ldots,i_k})_{i_1,\ldots,i_k\in\mathbb N}\text{ is summable in }A(X\lolli Y)\\
\mathrel{\stackrel{1}{\iff}}{}&\forall x\in AX,\forall y'\in AY^\perp,\scalar{\sum_{i_1,\ldots,i_k\in\mathbb N}\abs(f)_{i_1,\ldots,i_k}}{x\otimes y'}\text{ is defined and in }B\\
\mathrel{\stackrel{2}{\iff}}{}&\forall x\in AX,\forall y'\in AY^\perp,\scalar{\abs(h)}{x\otimes\Delta_{\mathbb N^k}\otimes y'}\text{ is defined and in }B\\
\mathrel{\stackrel{3}{\iff}}{}&\abs(h)\in A(X\otimes D^{\otimes k}\lolli Y)\iff h\in\Rmod{X\otimes D^{\otimes k}\lolli Y}.
\end{align*}
where (1) is the predual characterization of summability given in Proposition 33, (2) holds by the covering principle of Lemma 21 ($\{\Delta_{\mathbb N}\}$ is a covering), and (3) holds by the predual characterization of Theorem 23, using that $A(D^{\otimes k})=\{\Delta_{\mathbb N^k}\}^{\perp\perp}$ by Lemma 58.
\end{proof}
% SUPPORT PROOFS BEGIN -- official PDF Appendix A pp.29:20--21.
\section*{Original Appendix A: Proofs on web based models}
\paragraph{Proof of Lemma 20.} Observe that
\[\scalar{s\cdot x}{y}=\sum_{b\in|Y|}\left(\sum_{a\in|X|}s_{a,b}x_a\right)y_b\mathrel{\stackrel{1}{\dimp}}\sum_{b\in|Y|}\sum_{a\in|X|}(s_{a,b}x_ay_b)\mathrel{\stackrel{2}{\dequiv}}\sum_{a\in|X|,b\in|Y|}s_{a,b}x_ay_b=\scalar{s}{x\otimes y}\]
where (1) is by distributivity of multiplication over sums, and (2) is by (PA). Furthermore, if $s\cdot x$ is defined, then (1) turns into an equivalence $\dequiv$ (but not in general, consider $y_b=0$). The second item is proved in the same way.
\paragraph{Proof of Lemma 21.} Let $b_0\in|Y|$. We prove that $(s\cdot x)_{b_0}$ is defined. By covering, there exists $y\in F$ such that $y_{b_0}$ is invertible. By assumption, $(s_{a,b}x_ay_b)_{a\in|X|,b\in|Y|}$ is summable, so by Lemma 3 $(s_{a,b_0}x_ay_{b_0})_{a\in|X|}$ is summable. But then
\[y_{b_0}^{-1}\left(\sum_{a\in|X|}s_{a,b_0}x_ay_{b_0}\right)\dimp\sum_{a\in|X|}s_{a,b_0}x_ay_{b_0}y_{b_0}^{-1}=\sum_{a\in|X|}s_{a,b_0}x_a\]
so $(s_{a,b_0}x_a)_{a\in|X|}$ is summable.
\paragraph{Proof of Proposition 22.} We first prove (1) $\Leftrightarrow$ (2). By Remark 18,
\begin{align*}s\in A(X\lolli Y)&\iff\forall x\in AX,\forall y'\in AY,\scalar{s}{x\otimes y'}\text{ is defined and in }B,\\s^\perp\in A(Y^\perp\lolli X^\perp)&\iff\forall x\in AX,\forall y'\in AY,\scalar{s^\perp}{y'\otimes x}\text{ is defined and in }B.\end{align*}
We can easily check that $\scalar{s}{x\otimes y'}\dequiv\scalar{s^\perp}{y'\otimes x}$, which concludes the proof of equivalence. Next, we prove (1) $\Rightarrow$ (3). Let $s\in A(X\lolli Y)$ and $x\in AX$. First, observe that $s\cdot x$ is well-defined by covering principle (Lemma 21), since $AY^\perp$ is a covering. Then, for all $y'\in Y^\perp$, $\scalar{s\cdot x}{y'}\dequiv\scalar{s}{x\otimes y'}$ by Lemma 20. Since $s\perp x\otimes y'$ by assumption, we get that $s\cdot x\perp y'$. It concludes the proof of (3). The proof of the reverse implication (3) $\Rightarrow$ (1) relies on Lemma 20 in a similar same way. Finally, (2) $\Leftrightarrow$ (4) holds using (1) $\Leftrightarrow$ (3) on $s^\perp$.
\paragraph{Proof of Theorem 19.} Let $s\in A(X\lolli Y)$ and $t\in A(Y\lolli Z)$. First, let us prove that $t\cdot s$ is well-defined. By covering, for all $a_0\in|X|$ there exists $x\in AX$ such that $x_{a_0}$ is invertible. Then, $(t\cdot(s\cdot x))$ is defined by iterating Proposition 22 twice. But
\[(t\cdot(s\cdot x))_c=\sum_{b\in|Y|}t_{b,c}\left(\sum_{a\in|X|}s_{a,b}x_a\right)\dimp\sum_{b\in|Y|}\sum_{a\in|X|}t_{b,c}s_{a,b}x_a\mathrel{\stackrel{\mathrm{PA}}{\dequiv}}\sum_{a\in|X|,b\in|Y|}s_{a,b}t_{b,c}x_a.\]
It follows from (WPA) that $(s_{a_0,b}t_{b,x}x_{a_0})_{b\in|Y|}$ is summable, so by invertibility of $x_{a_0}$, $(s_{a_0,b}t_{b,c})_{b\in|Y|}$ is summable. That is, $t\cdot s$ is well-defined. We now prove that it belongs to $A(X\lolli Z)$ using Proposition 22. Let $x\in AX$. Then,
\begin{align*}((t\cdot s)\cdot x)_c&=\sum_{a\in|X|}\left(\sum_{b\in|Y|}s_{a,b}t_{b,c}\right)x_a\mathrel{\stackrel{1}{\dequiv}}\sum_{a\in|X|}\left(\sum_{b\in|Y|}s_{a,b}t_{b,c}x_a\right),\\(t\cdot(s\cdot x))_c&=\sum_{b\in|Y|}t_{b,c}\left(\sum_{a\in|X|}s_{a,b}x_a\right)\mathrel{\stackrel{2}{\dequiv}}\sum_{b\in|Y|}\left(\sum_{a\in|X|}s_{a,b}t_{b,c}x_a\right).\end{align*}
Be aware: (1) only holds because we already know that $t\cdot s$ is defined, and (2) only holds because we already know that $s\cdot x$ is defined. Then, by (PA), $(t\cdot s)\cdot x\dequiv t\cdot(s\cdot x)$. We conclude that $t\cdot s$ is well-defined and belongs to $A(X\lolli Z)$ by iterating Proposition 22 twice.
\paragraph{Proof of Theorem 23.} First, observe that (1) $\Rightarrow$ (2) is trivial.
\begin{description}
\item[(2) $\Rightarrow$ (3).] Let $s\in(P\otimes Q)^\perp$. Then for all $x\in P$, $s\cdot x$ is well-defined by covering principle (Lemma 21). Furthermore, for all $y'\in Q$, $\scalar{s}{x\otimes y'}\dequiv\scalar{s\cdot x}{y'}$ by Lemma 20. But $s\perp x\otimes y'$ so we deduce that $s\cdot x\perp y'$. Thus, $s\cdot x\in Q^\perp=AY$.
\item[(3) $\Rightarrow$ (2).] Assume that $s$ satisfies (3). Then we prove that $s\in(P\otimes AY^\perp)^\perp\subseteq(P\otimes Q)^\perp$. Let $x\in P$ and $y'\in(AY)^\perp$. Then by assumption $s\cdot x$ is defined so by Lemma 20 $\scalar{s}{x\otimes y'}\dequiv\scalar{s\cdot x}{y'}$ and we have that $s\perp x\otimes y'$. Thus, $s\in(P\otimes AY^\perp)^\perp$.
\item[(2) $\Leftrightarrow$ (4).] The equivalence holds by a similar argument, since $\scalar{s}{x\otimes y'}\dequiv\scalar{s^\perp}{y'\otimes x}$.
\item[(3) $\wedge$ (4) $\Rightarrow$ (1).] Assume that $s$ satisfies (3) and (4). First, we know from the proof of (3) $\Rightarrow$ (2) that $s\in(P\otimes AY^\perp)^\perp$. Second, by the implication (2) $\Rightarrow$ (4) applied on $P$ and $AY^\perp$, we know that for all $y'\in AY^\perp$, $s^\perp\cdot y'$ is defined and belongs to $AX^\perp$. So by the characterization of Proposition 22, we conclude that $s\in A(X\lolli Y)$.
\end{description}
\section*{Bibliographic context}
[22] Peter Hines. A categorical analogue of the monoid semiring construction. Mathematical Structures in Computer Science, 23(1):55--94, 2013. doi:10.1017/S0960129512000059.\par
[38] Paul-Andr\'e Melli\`es. Categorical semantics of linear logic. Panoramas et Synth\`eses, 27:1--196, 2009.\par
The complete published bibliography, examples and linear-logical structure Figure 1 remain in the accompanying original CC BY 4.0 published PDF, pp.29:8 and29:18--20. They are not additional counted targets.
